
This work examines the applicability of formal repair methods to LLM-generated code by measuring empirical preconditions before applying each method. The investigation integrates three repair approaches -- SynCode grammar-constrained decoding, Z3 SMT solving, and mypy type checking -- into a single Python-native pipeline and evaluates them on CodeLlama-7B output for HumanEval problems.

The central empirical finding is that across 2,680 completions from 134 HumanEval problems, CodeLlama-7B produces zero type errors. This means that mypy-based iterative repair, a technique used in LLM code improvement pipelines, cannot activate for this model on this benchmark. The precondition required by the type-checking feedback loop -- the presence of type errors -- is not satisfied.

\textbf{Motivation.} LLMs generate code with various error types: syntactic invalidity, type violations, arithmetic constraint failures, and logical errors. Different repair methods target different error types. Grammar-constrained decoding targets syntax by masking invalid tokens at generation time. SMT solving targets arithmetic constraint failures via symbolic reasoning. Type checking targets structural type violations via static analysis. Each method has a measurable precondition: the corresponding error type must actually appear in the generated code. If the precondition is not met, the method cannot activate regardless of its theoretical soundness.

\textbf{Gap.} Prior work on formal repair of LLM-generated code has not systematically measured the Failure Mode Distribution (FMD) -- the empirical breakdown of which failure types actually occur -- before applying formal repair methods. Repair method proposals are typically validated on end-to-end pass@k improvements without decomposing which failures the method can and cannot address. This gap has practical consequences: for instance, implementing mypy-based repair for CodeLlama-7B output would result in a repair loop that activates 0\% of the time, a fact that is not predictable from theoretical analysis alone.

\textbf{Approach.} This work takes a precondition-first view. Before asking which repair method achieves the highest pass@1, it asks which failure modes are present in a given setting and which repair channels can therefore activate. This is operationalized via an FMD pipeline: for each generated code sample, failures are classified as syntactic (via \texttt{ast.parse}), type-structural (via \texttt{mypy.api.run}), or functional/arithmetic (via test execution and Z3 eligibility scanning). Repair channels are gated on FMD-conditioned eligibility.

\textbf{Summary of findings.} The FMD measurement on CodeLlama-7B with HumanEval reveals: 97.5\% of failures are syntax-dominated; the type stratum contains zero problems across 134 problems and 2,680 samples; Z3 encoding is applicable to 33\% of problems (54/164) under integer-equality assertion scanning. SynCode directionally reduces AST failures (delta_ast = 0.075, consistent across two independent experiments). The mypy repair loop activates zero times. The SynCode-Z3 complementarity experiment (h-m3) was not executed.

\textbf{Contributions.}

\begin{enumerate}
\item A unified Python-native pipeline integrating SynCode, Z3-solver, and mypy with CodeLlama-7B on HumanEval without Docker dependencies.
\item An empirical FMD characterization of CodeLlama-7B on HumanEval: 97.5\% syntax-dominated, 0\% type stratum, 33\% Z3-eligible.
\item An empirical demonstration that mypy-based iterative repair is inapplicable to annotation-free code generation by CodeLlama-7B on HumanEval.
\item Implementation and validation of all components required for SynCode-Z3 complementarity testing.
\end{enumerate}
